% Source: https://artofproblemsolving.com/community/c3h1982025_mc_problem
% Stable identity: aops-topic-1982025-q01
A circle of radius $2$ through two points at distance $d$ exists only if $d\le4$; there are $2$ such circles if $d<4$ and $1$ (centred at the midpoint) if $d=4$. The $15$ pairs of points are:
\begin{itemize}
\item $4$ horizontal pairs at distance $2$ and $3$ vertical pairs at distance $2$: $2$ circles each, $14$ circles;
\item $4$ diagonal pairs at distance $2\sqrt2$: $2$ circles each, $8$ circles;
\item $2$ horizontal pairs at distance $4$, $(-1,\pm1)$ with $(3,\pm1)$: $1$ circle each, $2$ circles;
\item the remaining $2$ pairs, at distance $\sqrt{20}>4$: none.
\end{itemize}
That is $24$ (pair, circle) incidences. A circle through three of the points is counted three times. The circle centred at $(1,1)$ passes through $(-1,1)$, $(3,1)$ and $(1,-1)$, and the circle centred at $(1,-1)$ through $(-1,-1)$, $(3,-1)$ and $(1,1)$ (dashed); every other centre found above is at distance $2$ from only the two points that produced it (checking the $20$ centres directly). So there are $24-2\cdot3=18$ incidences belonging to circles through exactly two points, and each such circle is counted once.

\begin{asy}
usepackage("amsmath");
size(6cm);
defaultpen(fontsize(10pt));
pair[] p={(-1,1),(1,1),(3,1),(-1,-1),(1,-1),(3,-1)};
draw((-3.6,0)--(5.6,0),gray(0.7));draw((0,-3.4)--(0,3.4),gray(0.7));
draw(circle((1,1),2),dashed+gray(0.5));draw(circle((1,-1),2),dashed+gray(0.5));
draw(circle((0,1+sqrt(3)),2),linewidth(0.9));
for(int i=0;i<6;++i) dot(p[i]);
label("$(-1,1)$",p[0],SW,fontsize(8pt));label("$(1,1)$",p[1],S,fontsize(8pt));label("$(3,1)$",p[2],SE,fontsize(8pt));
label("$(-1,-1)$",p[3],SW,fontsize(8pt));label("$(1,-1)$",p[4],S,fontsize(8pt));label("$(3,-1)$",p[5],SE,fontsize(8pt));
label("\mbox{}",(min(currentpicture,user=true)+max(currentpicture,user=true))/2,black); // keeps the text after the figure black when LaTeX typesets the labels inline
\end{asy}

The answer is $\boxed{18}$. (One of them, centred at $(0,1+\sqrt3)$ through $(-1,1)$ and $(1,1)$, is drawn solid.)
